% Sign Families 328-372
\section{Sign Families: Basic Strokes and Numerals (M328--M372)}

\subsection{Vertical Stroke Signs (M328--M345)}

The vertical stroke series forms the core counting system:

\subsubsection{Single Strokes (M328--M335)}
\begin{itemize}
    \item M328: Single vertical stroke (unit marker)
    \item M329: Double vertical stroke (cardinal 2)
    \item M330: Triple vertical stroke (cardinal 3)
    \item M331--M335: Higher cardinal values
\end{itemize}

\textbf{Operational Function}: These signs implement cardinal arithmetic, representing discrete quantities in a unary counting system.

\subsection{Jar Signs (M346--M360)}

The jar series represents a class of container objects:

\subsubsection{Basic Jar Forms}
\begin{itemize}
    \item M346: Simple jar
    \item M347--M355: Jar variants (rim types, handle positions)
    \item M356--M360: Specialized containers
\end{itemize}

\textbf{Operational Function}: Instantiation operators introducing specific container types into the operational context.

\subsection{Fish Signs (M361--M372)}

The fish series forms a distinctive semantic-operational class:

\subsubsection{Major Fish Types}
\begin{itemize}
    \item M361: Basic fish
    \item M362--M368: Fish variants (orientation, species markers)
    \item M369--M372: Combined fish forms
\end{itemize}

\textbf{Operational Function}: These signs often appear in initial or near-initial positions, serving as domain specifiers or institutional markers.

\subsection{Distribution Patterns}

Analysis of co-occurrence:
\begin{itemize}
    \item Stroke signs rarely appear with other stroke signs (arithmetic closure)
    \item Jar signs co-occur with stroke signs (quantity specification)
    \item Fish signs often precede jar-stroke combinations (domain specification)
\end{itemize}
